\capituloPregunta{¿Cómo se construye un experimento factorial \(2^2\)?}

\section{Motivación}

Un diseño factorial \(2^2\) permite estudiar dos factores, cada uno con dos niveles. Con cuatro combinaciones se pueden estimar efectos principales e interacción.

\section{Tres historias para aprender a diseñar}

El material del curso permite contrastar tres situaciones:

\begin{enumerate}
  \item \textbf{Hongos en tortillas}: observación biológica interesante, útil para discutir contaminación, tiempo y condiciones ambientales, pero insuficiente como diseño si no se definen niveles, repeticiones y variable respuesta.
  \item \textbf{Lentejas}: sistema controlable para discutir tratamientos, blanco, germinación y presencia de hongos.
  \item \textbf{Cúrcuma--cloro en post-it}: práctica factorial \(2^2\) completa con blanco, seis repeticiones y medición con luxómetro.
\end{enumerate}

\begin{idea}
La diferencia entre observar y diseñar está en la estructura: factores, niveles, tratamientos, unidad experimental, controles, repeticiones y respuesta.
\end{idea}

\section{Caso central: cúrcuma y cloro en papel}

La práctica utiliza dos factores:

\begin{table}[H]
\centering
\begin{tabular}{ccc}
\toprule
Factor & Nivel bajo & Nivel alto\\
\midrule
Cúrcuma & 1 g & 4 g\\
Cloro comercial & 5 g & 10 g\\
\bottomrule
\end{tabular}
\caption{Factores y niveles del diseño \(2^2\).}
\end{table}

La unidad experimental es un post-it individual. Se aplica una solución con base común de agua y sal. Se mantienen constantes la distancia de aspersión, el número de aspersiones, el tipo de papel y el tiempo de secado.

\section{Tratamientos}

\begin{table}[H]
\centering
\begin{tabular}{cccc}
\toprule
Tratamiento & Cúrcuma & Cloro & Interpretación\\
\midrule
T1 & 1 g & 5 g & nivel bajo en ambos factores\\
T2 & 1 g & 10 g & cúrcuma baja, cloro alto\\
T3 & 4 g & 5 g & cúrcuma alta, cloro bajo\\
T4 & 4 g & 10 g & nivel alto en ambos factores\\
B1 & 0 g & 0 g & blanco experimental\\
\bottomrule
\end{tabular}
\caption{Tratamientos factoriales y blanco experimental.}
\end{table}

\section{Variable respuesta}

La variable cuantitativa principal es la transmisión o bloqueo de luz medido con luxómetro. También se pueden registrar variables visuales:

\begin{itemize}
  \item intensidad de color;
  \item uniformidad;
  \item contraste entre zona cubierta y expuesta;
  \item atractivo visual.
\end{itemize}

\section{Modelo}

El modelo factorial puede escribirse como:

\[
Y_{ijk} = \mu + A_i + B_j + (AB)_{ij} + \varepsilon_{ijk}
\]

donde \(A_i\) representa el efecto de la cúrcuma, \(B_j\) el efecto del cloro y \((AB)_{ij}\) la interacción.

\section{Error didáctico deliberado}

\begin{cuidado}
La práctica de hongos en tortillas no debe presentarse automáticamente como factorial \(2^2\). Primero debe preguntarse: ¿cuáles son los dos factores?, ¿cuáles son los dos niveles?, ¿cuál es la unidad experimental?, ¿cuántas repeticiones independientes existen?, ¿qué respuesta se medirá?
\end{cuidado}

\section{Actividad}

Transforme la observación de hongos en tortillas en un posible diseño experimental. Defina dos factores, dos niveles, una variable respuesta y un blanco o control. Después compare su diseño con el caso cúrcuma--cloro.

\section{Conexión}

Cuando el diseño sale del aula y entra a procesos manuales o industriales, aparecen restricciones: datos no balanceados, variabilidad operativa y límites materiales.
